\chapter{Le Processus d'Audit (ISO 19011)}
\label{chap:procaudita}

La norme \textbf{ISO 19011} fournit les lignes directrices pour auditer les systèmes de management. C'est la "bible" de l'auditeur, décrivant étape par étape comment mener une mission.

\section{Présentation de l'ISO 19011}

\subsection{Définition et Champ d'Application}

\begin{boxdef}
\textbf{ISO 19011} : Norme internationale intitulée "Lignes directrices pour l'audit des systèmes de management". Elle s'applique à tous les types d'audits (Qualité ISO 9001, Sécurité ISO 27001, Environnement ISO 14001, etc.).
\end{boxdef}

Elle définit les principes d'audit (Intégrité, Présentation équitable, Devoir de vigilance, Confidentialité, Indépendance, Approche fondée sur la preuve) et le processus complet.

\subsection{Vue d'ensemble du Processus}

Le processus d'audit est cyclique. La norme (Section 6) le divise en plusieurs phases séquentielles :
\begin{enumerate}
    \item Déclenchement de l'audit.
    \item Préparation des activités d'audit.
    \item Réalisation des activités d'audit.
    \item Préparation et distribution du rapport.
    \item Clôture et suivi.
\end{enumerate}

% ------------------------------------------------------------------------------
\section{Phase 1 : Déclenchement de l'Audit (ISO 19011, Section 6.2)}
% ------------------------------------------------------------------------------

Cette phase est administrative mais cruciale pour cadrer la mission.

\subsection{Établissement du contact initial}

L'auditeur rencontre le commanditaire pour définir les objectifs.

\begin{boxlizbiz}
\textbf{LiZbiZ :} Le commanditaire est le PDG. L'objectif n'est pas de "piéger la DSI", mais de "mesurer l'écart par rapport à l'ISO 27001" pour préparer la certification.
\end{boxlizbiz}

\subsection{Détermination de la faisabilité}

L'auditeur vérifie s'il a les ressources nécessaires :
\begin{itemize}
    \item Accès aux documents ?
    \item Disponibilité du personnel clé ?
    \item Temps imparti suffisant ?
\end{itemize}

\subsection{Constitution de l'équipe d'audit}

Pour LiZbiZ, l'équipe sera constituée d'un \textbf{Auditeur Principal} (l'étudiant chef de projet) et d'\textbf{Auditeurs Techniques} (les autres étudiants).

% ------------------------------------------------------------------------------
\section{Phase 2 : Préparation de l'Audit (ISO 19011, Section 6.3)}
% ------------------------------------------------------------------------------

C'est le travail "de bureau" avant d'aller sur le terrain.

\subsection{Revue documentaire (Pré-audit)}

L'auditeur lit les documents existants pour préparer sa grille de vérification.

\begin{boxlizbiz}
\textbf{LiZbiZ :} Vous demandez à consulter la "Politique de Sécurité" et "La procédure de gestion des accès". Si ces documents n'existent pas, c'est déjà une Non-Conformité potentielle identifiée avant même de commencer.
\end{boxlizbiz}

\subsection{Plan d'Audit (Audit Plan)}

C'est le document contractuel validé par l'audité. Il précise le "Qui, Quoi, Où, Quand".

\begin{table}[htbp]
\centering
\small
\begin{tabular}{|p{3cm}|p{3cm}|p{3cm}|p{3cm}|}
    \hline
    \textbf{Date / Heure} & \textbf{Activité} & \textbf{Auditeur} & \textbf{Audité} \\
    \hline
    09h00 - 09h30 & Réunion d'ouverture & Chef d'équipe & DSI + Equipe \\
    \hline
    09h30 - 11h00 & Audit technique (Pare-feu, Logs) & Auditeur Tech & Admin Sys \\
    \hline
    11h00 - 12h30 & Audit organisationnel (RH) & Auditeur Orga & DRH \\
    \hline
    14h00 - 15h00 & Synthèse équipe & Équipe complète & - \\
    \hline
    15h00 - 16h00 & Réunion de clôture & Chef d'équipe & DSI + PDG \\
    \hline
\end{tabular}
\caption{Extrait du Plan d'Audit LiZbiZ}
\end{table}

\subsection{Feuille de Route (Checklist)}

C'est l'outil de travail de l'auditeur sur le terrain. Elle liste les questions à poser et les points à vérifier.

\begin{boxdef}
\textbf{Checklist (Feuille de route) :} Document détaillé traduisant les exigences de la norme (ex: ISO 27001) en questions fermées ou ouvertes.
\end{boxdef}

Exemple de checklist pour l'ISO 27001 Annexe A.9 (Contrôle d'accès) :
\begin{itemize}
    \item \textit{Question :} "Montrez-moi la liste des utilisateurs ayant des privilèges administrateurs."
    \item \textit{Question :} "Comment procédez-vous lors du départ d'un collaborateur ?"
\end{itemize}

% ------------------------------------------------------------------------------
\section{Phase 3 : Réalisation de l'Audit (ISO 19011, Section 6.4)}
% ------------------------------------------------------------------------------

C'est le moment de la vérité, sur le terrain (ou dans le lab GNS3).

\subsection{Réunion d'Ouverture (Opening Meeting)}

C'est une réunion formelle pour confirmer le plan.
\textbf{Ordre du jour type :}
\begin{itemize}
    \item Présentation des participants.
    \item Rappel du périmètre et des objectifs.
    \item Confirmation des ressources et accès.
    \item Rappel des règles de confidentialité.
\end{itemize}

\subsection{Collecte et Vérification des Informations}

L'auditeur collecte des preuves par trois moyens (Technique des 3 "E") :
\begin{enumerate}
    \item \textbf{Entretiens :} Questionner le personnel.
    \item \textbf{Examen documentaire :} Lire les procédures, regarder les logs.
    \item \textbf{Observation :} Regarder les activités réelles (ex: voir si l'écran est verrouillé).
\end{enumerate}

\begin{boxalert}
\textbf{Règle de l'échantillonnage :} L'auditeur ne peut pas tout vérifier. Il sélectionne des échantillons (ex: 5 comptes utilisateurs au hasard sur 100). Si 2 sur 5 sont mal gérés, il conclura à une anomalie généralisée.
\end{boxalert}

\subsection{Réunion de Clôture (Closing Meeting)}

À la fin de l'audit, l'équipe présente ses conclusions à l'audité.
\begin{itemize}
    \item C'est un moment de communication délicate.
    \item On présente les constats (faits objectifs).
    \item On annonce les Non-Conformités (NC).
    \item On ne donne pas encore les recommandations (ce n'est pas le rôle de l'auditeur de faire le travail de l'audité, sauf audit interne conseil).
\end{itemize}

% ------------------------------------------------------------------------------
\section{Phase 4 et 5 : Rapport et Suivi (ISO 19011, Sections 6.5 à 6.7)}
% ------------------------------------------------------------------------------

\subsection{Le Rapport d'Audit}

Il doit être clair, complet et factuel. Il contient :
\begin{itemize}
    \item Le périmètre et les objectifs.
    \item Le déroulement de l'audit.
    \item Les conclusions (Forces, Faiblesses).
    \item Les Non-Conformités détaillées (Exigence non satisfaite + Preuve).
\end{itemize}

\subsection{Le Suivi (Follow-up)}

L'audit ne s'arrête pas au rapport. L'audité doit mettre en place des actions correctives. L'auditeur doit vérifier l'efficacité de ces actions (souvent lors d'un audit de suivi ultérieur).

% ------------------------------------------------------------------------------
\section{Cas LiZbiZ : Simulation des Réunions}
% ------------------------------------------------------------------------------

\subsection{Exercice de Simulation (Jeu de rôle)}

\begin{boxlizbiz}
\textbf{Mise en situation :}
Les étudiants se répartissent les rôles : Auditeurs vs Audités (DSI, Admin Sys, RH).

\textbf{Scénario Réunion d'Ouverture :}
Les auditeurs doivent présenter le plan d'audit. L'Admin Sys (joué par un étudiant) tente de refuser l'accès aux logs en disant "C'est trop sensible, vous allez casser le système".
\textbf{Réponse attendue :} L'auditeur doit rappeler le mandat du PDG et les règles de l'audit (non-modification du système).
\end{boxlizbiz}

% ==============================================================================
% Fichier : 03_module_audit.tex (Extrait Chapitre 2)
% Description : Le Processus d'Audit selon ISO 19011
% ==============================================================================